\section{確率的プログラミング言語}
\subsection{授業内容}
\subsubsection{科目の中でのこのコマの位置づけ}
これまではRにビルトインされている確率関数を使って，シミュレーションや検定のロジックを確認してきた。ここまでは，確率を反復試行におけるある実現値の出現割合と考えてきたことになる。この考えに基づく確率の定義は頻度主義的ともいわれるが，同じ確率を使った推論でも他の方法が存在する。ベイズ的確率の考えがそれで，この場合の推論は事前分布と確率モデルから導出される事後分布を活用して行われる。
心理統計でよく用いられるモデルは，一様分布を事前分布とし，確率モデルは正規分布を仮定することがほとんどであり，事後分布も解析的に導出することができるが，ここではより一般的な確率モデルに応用することを考える。
確率的プログラミング言語を持ちることで，尤度と事前分布を設定することで事後分布発生器を作ることができ，いかなる事後分布からでも乱数を生成することができる。乱数によって確率分布が近似できることはすでに学んだとおりである。
ここでは専門の確率的プログラミング言語Stanを導入し，以後の学習に備える。
\subsubsection{コマ主題細目}
\begin{description}
	\item[ベイズ推論の基礎] 推定法の一種，ベイズ推定はベイズの定理に基づいている。改めてベイズの定理を解説し，条件付き確率や尤度，事前分布，事後分布といった用語についても改めて確認する。
	\begin{flushright}
		$\to$\textcite{Maeda2017}のPp104--123ほか枚挙に遑がない
	\end{flushright}
	\item[事後分布の生成] マルコフ連鎖モンテカルロ法は，事後分布の生成と乱数の生成が合体した計算技術である。事後分布の関数の形はわからなくとも，そこから乱数を生成することで形状を近似し，分析することができる。これらの基本的な仕組みを理解する。
	\begin{flushright}
		$\to$\textcite{Maeda2017}のPp.147--194
	\end{flushright}
	\item[確率的プログラミング言語Stan] マルコフ連鎖モンテカルロ法の演習にあたって，確率的プログラミング言語Stanを導入する。Stanは変数の宣言が必要であること，ブロックに分割されていること，行の終わりにセミコロンを入れることなど，R言語よりもC言語に近い書き方が必要である。RStudioにはStanファイルのサンプルも入っているので，これらを活用して簡単なコードを書いてみる。
	\begin{flushright}
		$\to$\textcite{Maeda2017}のPp.407--425.
	\end{flushright}
	\item[事後分布による解析] StanをRから利用すると，結果として得られるオブジェクトに非常に多くの情報を含んでいることになる。ここではcmdstanrの出力結果を確認し，またコンパクトに出力をまとめる関数を作りながら，事後分布をどのように可視化，記述するかについて学ぶ。
\end{description}
\subsubsection{キーワード}
\begin{itemize}
	\item ベイズ推定
	\item 事前分布，事後分布
	\item 確率的プログラミング言語
	\item Stan
\end{itemize}
\subsection{授業情報}
\paragraph{コマの展開方法}
講義/遠隔/演習
\subsubsection{予習・復習課題}
\paragraph{予習} Stanの導入は大学のPC教室を利用するが，そうでなければローカルに環境を構築する必要がある。ローカル環境での利用を希望するものは，事前に調べて用意しておく。
\paragraph{復習} Stanのサンプルコード，可視化の関数やパッケージなど，演習の分量が多いため，時間内に終わらなかったものは資料を参考に必ずフォローアップしておくこと。